\section{Tarefa 6.1: Projeto e simula\c{c}\~ao de um circuito multiplicador de 2 bits}

O projeto foi feito a partir da gera\c{c}\~ao dos produtos parciais com portas AND entre os bits de entrada $A_1A_0$ e $B_1B_0$. O bit menos significativo \'e dado diretamente por $C_0 = A_0B_0$. Para o segundo bit, somam-se os termos cruzados $A_0B_1$ e $A_1B_0$ com uma porta XOR, produzindo $M_1$. O terceiro bit \'e obtido com a soma (XOR) entre $A_1B_1$ e o carry da soma anterior, que \'e $(A_0B_1)(A_1B_0)$, resultando em $M_2$. Por fim, o bit mais significativo $M_3$ corresponde ao carry final, implementado no esquem\'atico como o produto $A_1A_0B_1B_0$.

\begin{figure}[H]
    \centering
    \includegraphics[width=0.95\textwidth]{esquematico.png}
    \caption{Diagrama esquem\'atico do multiplicador de 2 bits.}
\end{figure}

\subsection{Equa\c{c}\~oes do circuito}

\[
    P_0 = C_0 = A_0 \cdot B_0
\]

\[
    P_1 = M_1 = (A_0 \cdot B_1) \oplus (A_1 \cdot B_0)
\]

\[
    P_2 = M_2 = (A_1 \cdot B_1) \oplus \left[(A_0 \cdot B_1) \cdot (A_1 \cdot B_0)\right]
\]

\[
    P_3 = M_3 = A_1 \cdot A_0 \cdot B_1 \cdot B_0
\]

\subsection{Simula\c{c}\~ao funcional}


\begin{figure}[H]
    \centering
    \includegraphics[width=0.9\textwidth]{\detokenize{produto 1011.png}}
    \caption{Temporiza\c{c}\~ao para $(10)_b \times (11)_b$ (entrada 1011, sa\'ida 0110).}
\end{figure}

\begin{figure}[H]
    \centering
    \includegraphics[width=0.9\textwidth]{\detokenize{produto 1000.png}}
    \caption{Temporiza\c{c}\~ao para $(10)_b \times (00)_b$ (entrada 1000, sa\'ida 0000).}
\end{figure}

\begin{figure}[H]
    \centering
    \includegraphics[width=0.9\textwidth]{\detokenize{produto 0101.png}}
    \caption{Temporiza\c{c}\~ao para $(01)_b \times (01)_b$ (entrada 0101, sa\'ida 0001).}
\end{figure}

\begin{figure}[H]
    \centering
    \includegraphics[width=0.9\textwidth]{\detokenize{produto 1111.png}}
    \caption{Temporiza\c{c}\~ao para $(11)_b \times (11)_b$ (entrada 1111, sa\'ida 1001).}
\end{figure}

\begin{table}[H]
    \centering
    \caption{Resultados esperados para os casos simulados.}
    \begin{tabular}{ccc}
        \toprule
        Opera\c{c}\~ao          & Resultado esperado & Resultado simulado \\
        \midrule
        $(10)_b \times (11)_b$ & $(0110)_b$         & $(0110)_b$         \\
        $(10)_b \times (00)_b$ & $(0000)_b$         & $(0000)_b$         \\
        $(01)_b \times (01)_b$ & $(0001)_b$         & $(0001)_b$         \\
        $(11)_b \times (11)_b$ & $(1001)_b$         & $(1001)_b$         \\
        \bottomrule
    \end{tabular}
\end{table}

\subsection{Coment\'arios}

    Os resultados simulados coincidem com os valores esperados para os quatro casos pedidos no roteiro. Isso confirma que o esquem\'atico implementa corretamente a multiplica\c{c}\~ao de 2 bits, com sa\'ida organizada como $P_3P_2P_1P_0 = M_3M_2M_1C_0$.
